\documentclass[10pt]{article}
\usepackage[utf8]{inputenc}
\usepackage[vietnamese]{babel}
\usepackage{amsmath,amssymb,amsfonts}
\usepackage{newtxtext,newtxmath}
\usepackage{geometry}
\usepackage{xcolor}
\usepackage{tikz}
\usepackage{fancyhdr}

\geometry{
  paperwidth=210mm,
  paperheight=280mm,
  left=16mm,
  right=16mm,
  top=14mm,
  bottom=14mm
}

\definecolor{stewartcyan}{RGB}{0, 130, 165}

\pagestyle{fancy}
\fancyhf{}
\renewcommand{\headrulewidth}{0pt}
\fancyhead[R]{{\fontsize{8.5pt}{10pt}\selectfont\sffamily\textbf{\color{stewartcyan}MỤC 4.1}\quad Giá trị lớn nhất và giá trị nhỏ nhất}\qquad{\fontsize{9.5pt}{11pt}\selectfont\sffamily\textbf{287}}}

\newcommand{\graphicon}{%
  \begin{tikzpicture}[baseline=-0.6ex, scale=0.85]
    \draw[color=stewartcyan, line width=0.6pt, rounded corners=0.5pt] (0,0) rectangle (9pt,9pt);
    \draw[color=stewartcyan, line width=0.35pt] (0.8pt, 4.5pt) -- (8.2pt, 4.5pt);
    \draw[color=stewartcyan, line width=0.35pt] (4.5pt, 0.8pt) -- (4.5pt, 8.2pt);
    \draw[color=stewartcyan, line width=0.7pt, smooth] plot coordinates {(1.2pt, 2pt) (3.2pt, 7.2pt) (5.8pt, 1.8pt) (7.8pt, 7pt)};
  \end{tikzpicture}%
}

\setlength{\parindent}{0pt}
\setlength{\parskip}{0pt}

\begin{document}
\fontsize{9pt}{11.3pt}\selectfont

\noindent
\begin{minipage}[t]{0.485\textwidth}
\makebox[1.6em][l]{\textbf{12.}}\textbf{(a)} Phác thảo đồ thị của một hàm số trên $[-1, 2]$ có giá trị lớn nhất tuyệt đối nhưng không có cực đại địa phương.\\[1pt]
\makebox[1.6em][l]{}\textbf{(b)} Phác thảo đồ thị của một hàm số trên $[-1, 2]$ có cực đại địa phương nhưng không có giá trị lớn nhất tuyệt đối.

\vspace{6pt}
\makebox[1.6em][l]{\textbf{13.}}\textbf{(a)} Phác thảo đồ thị của một hàm số trên $[-1, 2]$ có giá trị lớn nhất tuyệt đối nhưng không có giá trị nhỏ nhất tuyệt đối.\\[1pt]
\makebox[1.6em][l]{}\textbf{(b)} Phác thảo đồ thị của một hàm số trên $[-1, 2]$ gián đoạn nhưng có cả giá trị lớn nhất tuyệt đối và giá trị nhỏ nhất tuyệt đối.

\vspace{6pt}
\makebox[1.6em][l]{\textbf{14.}}\textbf{(a)} Phác thảo đồ thị của một hàm số có hai cực đại địa phương, một cực tiểu địa phương, và không có giá trị nhỏ nhất tuyệt đối.\\[1pt]
\makebox[1.6em][l]{}\textbf{(b)} Phác thảo đồ thị của một hàm số có ba cực tiểu địa phương, hai cực đại địa phương, và bảy điểm tới hạn.

\vspace{6pt}
{\color{stewartcyan}\textbf{15--28}}\quad Phác thảo đồ thị của $f$ bằng tay và sử dụng đồ thị phác thảo đó để tìm các giá trị lớn nhất và nhỏ nhất tuyệt đối cũng như địa phương của $f$. (Sử dụng các đồ thị và phép biến đổi đồ thị ở Mục 1.2 và 1.3.)

\vspace{5pt}
\makebox[1.6em][l]{\textbf{15.}}$f(x) = 3 - 2x, \quad x \ge -1$\\[3pt]
\makebox[1.6em][l]{\textbf{16.}}$f(x) = x^2, \quad -1 \le x < 2$\\[3pt]
\makebox[1.6em][l]{\textbf{17.}}$f(x) = 1/x, \quad x \ge 1$\\[3pt]
\makebox[1.6em][l]{\textbf{18.}}$f(x) = 1/x, \quad 1 < x < 3$\\[3pt]
\makebox[1.6em][l]{\textbf{19.}}$f(x) = \sin x, \quad 0 \le x < \pi/2$\\[3pt]
\makebox[1.6em][l]{\textbf{20.}}$f(x) = \sin x, \quad 0 < x \le \pi/2$\\[3pt]
\makebox[1.6em][l]{\textbf{21.}}$f(x) = \sin x, \quad -\pi/2 \le x \le \pi/2$\\[3pt]
\makebox[1.6em][l]{\textbf{22.}}$f(t) = \cos t, \quad -3\pi/2 \le t \le 3\pi/2$\\[3pt]
\makebox[0.56\linewidth][l]{\makebox[1.6em][l]{\textbf{23.}}$f(x) = \ln x, \quad 0 < x \le 2$}\textbf{24.} $f(x) = |x|$\\[3pt]
\makebox[0.56\linewidth][l]{\makebox[1.6em][l]{\textbf{25.}}$f(x) = 1 - \sqrt{x}$}\textbf{26.} $f(x) = e^x$\\[3pt]
\makebox[1.6em][l]{\textbf{27.}}$f(x) = \begin{cases} x^2 & \text{nếu } -1 \le x \le 0 \\ 2 - 3x & \text{nếu } 0 < x \le 1 \end{cases}$\\[4pt]
\makebox[1.6em][l]{\textbf{28.}}$f(x) = \begin{cases} 2x + 1 & \text{nếu } 0 \le x < 1 \\ 4 - 2x & \text{nếu } 1 \le x \le 3 \end{cases}$

\vspace{8pt}
{\color{stewartcyan}\textbf{29--48}}\quad Tìm các điểm tới hạn của hàm số.

\vspace{5pt}
\begin{tabular}{@{}p{0.485\linewidth}@{\hspace{0.03\linewidth}}p{0.485\linewidth}@{}}
\textbf{29.} $f(x) = 3x^2 + x - 2$ & \textbf{30.} $g(v) = v^3 - 12v + 4$ \\[4pt]
\textbf{31.} $f(x) = 3x^4 + 8x^3 - 48x^2$ & \textbf{32.} $f(x) = 2x^3 + x^2 + 8x$ \\[4pt]
\textbf{33.} $g(t) = t^5 + 5t^3 + 50t$ & \textbf{34.} $A(x) = |3 - 2x|$ \\[4pt]
\textbf{35.} $g(y) = \dfrac{y - 1}{y^2 - y + 1}$ & \textbf{36.} $h(p) = \dfrac{p - 1}{p^2 + 4}$ \\[8pt]
\textbf{37.} $p(x) = \dfrac{x^2 + 2}{2x - 1}$ & \textbf{38.} $q(t) = \dfrac{t^2 + 9}{t^2 - 9}$ \\[8pt]
\textbf{39.} $h(t) = t^{3/4} - 2t^{1/4}$ & \textbf{40.} $g(x) = \sqrt[3]{4 - x^2}$ \\[4pt]
\textbf{41.} $F(x) = x^{4/5}(x - 4)^2$ & \textbf{42.} $h(x) = x^{-1/3}(x - 2)$ \\[4pt]
\textbf{43.} $f(x) = x^{1/3}(4 - x)^{2/3}$ & \textbf{44.} $f(\theta) = \theta + \sqrt{2}\cos\theta$
\end{tabular}
\end{minipage}%
\hfill
\begin{minipage}[t]{0.485\textwidth}
\begin{tabular}{@{}p{0.485\linewidth}@{\hspace{0.03\linewidth}}p{0.485\linewidth}@{}}
\textbf{45.} $f(\theta) = 2\cos\theta + \sin^2\theta$ & \textbf{46.} $p(t) = te^{4t}$ \\[4pt]
\textbf{47.} $g(x) = x^2\ln x$ & \textbf{48.} $B(u) = 4\tan^{-1}u - u$
\end{tabular}

\vspace{6pt}
\graphicon\ {\color{stewartcyan}\textbf{49--50}}\quad Một công thức cho \textit{đạo hàm} của hàm số $f$ được cho. Hàm số $f$ có bao nhiêu điểm tới hạn?

\vspace{4pt}
\begin{tabular}{@{}p{0.485\linewidth}@{\hspace{0.03\linewidth}}p{0.485\linewidth}@{}}
\textbf{49.} $f'(x) = 5e^{-0.1|x|}\sin x - 1$ & \textbf{50.} $f'(x) = \dfrac{100\cos^2 x}{10 + x^2} - 1$
\end{tabular}

\vspace{8pt}
{\color{stewartcyan}\textbf{51--66}}\quad Tìm giá trị lớn nhất tuyệt đối và giá trị nhỏ nhất tuyệt đối của $f$ trên đoạn đã cho.

\vspace{4pt}
\makebox[1.6em][l]{\textbf{51.}}$f(x) = 12 + 4x - x^2, \quad [0, 5]$\\[3pt]
\makebox[1.6em][l]{\textbf{52.}}$f(x) = 5 + 54x - 2x^3, \quad [0, 4]$\\[3pt]
\makebox[1.6em][l]{\textbf{53.}}$f(x) = 2x^3 - 3x^2 - 12x + 1, \quad [-2, 3]$\\[3pt]
\makebox[1.6em][l]{\textbf{54.}}$f(x) = x^3 - 6x^2 + 5, \quad [-3, 5]$\\[3pt]
\makebox[1.6em][l]{\textbf{55.}}$f(x) = 3x^4 - 4x^3 - 12x^2 + 1, \quad [-2, 3]$\\[3pt]
\makebox[1.6em][l]{\textbf{56.}}$f(t) = (t^2 - 4)^3, \quad [-2, 3]$\\[3pt]
\makebox[1.6em][l]{\textbf{57.}}$f(x) = x + \dfrac{1}{x}, \quad [0.2, 4]$\\[4pt]
\makebox[1.6em][l]{\textbf{58.}}$f(x) = \dfrac{x}{x^2 - x + 1}, \quad [0, 3]$\\[4pt]
\makebox[1.6em][l]{\textbf{59.}}$f(t) = t - \sqrt[3]{t}, \quad [-1, 4]$\\[3pt]
\makebox[1.6em][l]{\textbf{60.}}$f(x) = \dfrac{e^x}{1 + x^2}, \quad [0, 3]$\\[4pt]
\makebox[1.6em][l]{\textbf{61.}}$f(t) = 2\cos t + \sin 2t, \quad [0, \pi/2]$\\[3pt]
\makebox[1.6em][l]{\textbf{62.}}$f(\theta) = 1 + \cos^2\theta, \quad [\pi/4, \pi]$\\[3pt]
\makebox[1.6em][l]{\textbf{63.}}$f(x) = x^{-2}\ln x, \quad \left[\frac{1}{2}, 4\right]$\\[3pt]
\makebox[1.6em][l]{\textbf{64.}}$f(x) = xe^{x/2}, \quad [-3, 1]$\\[3pt]
\makebox[1.6em][l]{\textbf{65.}}$f(x) = \ln(x^2 + x + 1), \quad [-1, 1]$\\[3pt]
\makebox[1.6em][l]{\textbf{66.}}$f(x) = x - 2\tan^{-1}x, \quad [0, 4]$

\vspace{8pt}
\makebox[1.6em][l]{\textbf{67.}}Nếu $a$ và $b$ là các số dương, hãy tìm giá trị lớn nhất của $f(x) = x^a(1 - x)^b, \quad 0 \le x \le 1$.

\vspace{4pt}
\graphicon\ \textbf{68.}\quad Sử dụng đồ thị để ước lượng các điểm tới hạn của $f(x) = |1 + 5x - x^3|$ chính xác đến một chữ số thập phân.

\vspace{4pt}
\graphicon\ {\color{stewartcyan}\textbf{69--72}}\quad \textbf{(a)} Sử dụng đồ thị để ước lượng các giá trị lớn nhất và nhỏ nhất tuyệt đối của hàm số chính xác đến hai chữ số thập phân.\\[2pt]
\hspace*{1.8em}\textbf{(b)} Sử dụng phép tính vi tích phân để tìm các giá trị lớn nhất và nhỏ nhất tuyệt đối chính xác.

\vspace{4pt}
\makebox[1.6em][l]{\textbf{69.}}$f(x) = x^5 - x^3 + 2, \quad -1 \le x \le 1$\\[3pt]
\makebox[1.6em][l]{\textbf{70.}}$f(x) = e^x + e^{-2x}, \quad 0 \le x \le 1$\\[3pt]
\makebox[1.6em][l]{\textbf{71.}}$f(x) = x\sqrt{x - x^2}$\\[3pt]
\makebox[1.6em][l]{\textbf{72.}}$f(x) = x - 2\cos x, \quad -2 \le x \le 0$
\end{minipage}

\vfill
\begin{center}
\fontsize{6.5pt}{8.2pt}\selectfont\color{black!75}
Bản quyền 2021 Cengage Learning. Bảo lưu mọi quyền. Không được sao chép, quét hoặc nhân bản, toàn bộ hoặc một phần. WCN 02-200-203\\
Đánh giá biên tập cho thấy rằng mọi nội dung bị lược bỏ không ảnh hưởng đáng kể đến trải nghiệm học tập tổng thể. Cengage Learning có quyền gỡ bỏ nội dung bổ sung vào bất kỳ lúc nào nếu các hạn chế về quyền sau này yêu cầu.
\end{center}

\end{document}